\scanpage{166}
\uline{Axiomes supplémentaires} \uline{Associativité} de
l'homomorphisme $\Delta$\,: correspond à un axiome
d'associativité pour les produits tensoriels.
\margin{On exprime ici que les hom. tensoriels classiques sont des $\underline U$-homomorphismes}
% slanted

\uline{Commutativité} de $\Delta$\,: l'hom canonique
$V\underset{k}{\otimes}W\longrightarrow W\underset{k}{\otimes}V$ doit être compatible
avec les $\underline U$-structures. On ne
parlera par la suite que de bifoncteurs
tensoriels associatifs et commutatifs, i.e. $\Delta$
associatif et commutatif. --- \uline{Compatibilité}
\uline{avec l'augmentation} \inserted{(i.e. $\Delta(u)=1\otimes u+u\otimes1+{}$\unsure{él. mixtes|élém. mixtes} \unsure{si} $\unsure{\varepsilon(u)=0}$)} $\Longleftrightarrow$ Le produit tensoriel
d'un $\underline U$-module $V$ avec le $\underline U$-module
trivial $k$ est le \deleted{$\underline U$-module trivial} $V$

\uline{Compatibilité avec $\vee$}\,: \deleted{\unsure{Si} $k$ $V,W$} Si $V,W\in K^{\underline U}$,
alors $(V\underset{k}{\otimes}W)' \simeq V'\underset{k}{\otimes}W'$ --- \unsure{tant que|en tant que}
$k$-modules, on veut que ce soit vrai
en tant que $\underline U$-modules \unsure{permis|\deleted{permis}}. Ceci
signifie qu'on a commutativité dans
\diagram{\begin{tikzcd}[ampersand replacement=\&] U\arrow[r,"\Delta"]\arrow[d,"\vee"] \& U\otimes U\arrow[d,"\vee\otimes\vee"] \\ U\arrow[r,"\Delta"] \& U\otimes U \end{tikzcd}}
\diagram{\begin{tikzcd}[ampersand replacement=\&] \deleted{\unsure{U}}\arrow[r,dash]\arrow[d] \& \deleted{\unsure{U\otimes U}}\arrow[d] \\ \deleted{\unsure{U}}\arrow[r] \& \deleted{\unsure{U\otimes U}} \end{tikzcd}}
i.e. $\Delta\vee=\deleted{\unsure{\Delta}}(\vee\otimes\vee)\Delta$, i.e $\vee$ est un
isomorphisme de co-algèbre.

% archive sheet number: 165
